% Page 019 — page imprimée 15
% Une épidémie à Meyrueis en 1768
% Reconstruction éditoriale en LaTeX.

{\fontsize{10.25}{12.1}\selectfont

\sectiondate{Une épidémie à Meyrueis}{1768}


\begin{center}
\fbox{%
\begin{minipage}{0.91\textwidth}
\itshape
Ce compte rendu d’inspection médicale à Meyrueis et dans les environs a été rédigé en 1768, vingt ans avant la révolution. On peut le consulter au Service du Patrimoine de la Médiathèque de Montpellier sous la cote 19.180.

\vspace{0.35em}

Il m’a vivement intéressé pour trois raisons :

\vspace{0.15em}

1/ \quad La description assez précise de cette épidémie dénommée « peste blanche » conséquence du manque d’hygiène, de la pauvreté et d’une nourriture insuffisante et souvent avariée.

\vspace{0.15em}

2/ \quad L’évocation de la vie quotidienne des paysans de Meyrueis et des vallées des environs à la suite de mauvaises récoltes dues à des sécheresses persistante et à des hivers rigoureux. Les maisons ne sont que des masures sans aération où « l’on meurt de froid ou de chaud », la promiscuité est générale, excellent terreau pour la tuberculose et toutes les maladies infectieuses.

\vspace{0.15em}

3/ \quad Le style et les tournures surannées de ces lignes.

\vspace{0.15em}

Ce mémoire complète donc ce que l’on peut savoir de la vie de nos aïeux à Meyrueis mais aussi à Carnac près de Rousses où la famille Couderc résidait alors.
\end{minipage}}
\end{center}

\vspace{5mm}

\begin{center}
{\fontsize{18}{20}\selectfont\bfseries MEMOIRE\par}
\vspace{0.55em}
{\bfseries
sur la maladie épidémique\\
qui a régné à Meyrueis et ses environs
}
\vspace{0.15em}

Présenté à Monsieur de Saint Priest\\
Le 3 Septembre 1768

\vspace{0.35em}

{\bfseries Par Tandon, Docteur en Médecine\par}
Imprimé par Auguste François Rochard, Place du petit Scel\\
Montpellier 1769
\end{center}

\vspace{0.55em}

C’est une maladie qui a régné assez longtemps dans une partie très considérable des Hautes Cévennes, manifestée par des fièvres intermittentes d’un très mauvais caractère. Elles ont d’abord sévi à St Jean de Védas, la Verune, Fabrègues, Gigean, Loupian, Cournon venant des exhalaisons des étangs par fortes chaleurs et de mauvaises récoltes. Le plus grand nombre des habitants est obligé de se nourrir d’\textit{azéroles}\footnote{Fruit de l’azérolier, jaune ou rouge, ressemblant à une petite pomme}, de mûres de haies, de figues et de raisins qui n’étaient pas mûrs et de petites carpes, muges, anguilles et loups.

Puis l’épidémie s’est répandue sur Marsillargue, Clermont, Lodève, Mèze et enfin les Hautes Cévennes.

\vspace{0.3em}

}

{\fontsize{10.25}{12.1}\selectfont
La maladie épidémique qui règne depuis environ cinq mois, parmi les gens de la campagne à Meyrueis et ses environs, Lanuéjols, Montjardin, Serviglières et quelques autres lieux moins considérables, est une fièvre, maligne, bilieuse, qui, à ce qu’on affirme, a fait de grands ravages à Séverac en Rouergue et a parcouru le Vallon du Tarn. Cette maladie prélude pendant plusieurs jours des malaises, des lassitudes, des pesanteurs d’estomac, le dégoût et le cours de ventre bilieux... la langue auparavant rouge, sèche, se racornit, se fend en différents sens et se retire dans le fond de la bouche, ou se couvre d’une croûte noirâtre et raboteuse qui s’étend sur les gencives et les lèvres qu’elle fait paraître brûlées et dans cet état ils ne peuvent ni parler ni avaler à moins qu’on ait l’attention d’humecter souvent ces parties.

\begin{itemize}
\item Il est très douteux qu’on ait vus des bubons.
\item Le terme le plus ordinaire est du quatorzième ou quinzième jour, parfois 22 ou 23.
\item Les enfants en ont tous réchappé mais à Lanuéjols il y a eu 45 morts.
\end{itemize}

\medskip
\textit{Quelles sont les causes les plus apparentes de cette épidémie ?}


« L’excessive misère qui règne depuis longtemps parmi les habitants de ces montagnes. La mauvaise nourriture et la malpropreté qui en sont inséparables, les alternatives de froid et de chaud, les orages mais surtout les brouillards qu’ils y ont presque continuellement essuyés. Une longue suite de mauvaises récoltes a rendu cette misère si sensible qu’à ne considérer que l’état actuel de ces infortunés, on ne saurait se persuader qu’ils aient jamais pu être les pères nourriciers de la plus grande partie des Cévennes...

... relégués dans le fond des vallons, tous couverts de haillons et ayant à peine une chemise de rechange, ils y habitent de misérables chaumières placées sur le bord de quelque torrent et entourées de fumiers de buis qui exhalent une odeur des plus insupportable....

... obligé de pénétrer dans ces chaumières on est tout étonné de voir que l’air intérieur de ces réduits obscurs et malsains y est à grand peine renouvelé par l’ouverture étroite d’une ample cheminée, d’une porte écrasée et d’une fente qui tient lieu de fenêtre pratiquée sur l’un des côtés de cette porte et de n’y trouver que quelques mauvaises planches mal assurées dans le mur, dont ils n’ont maintenant que faire, mais qui, autrefois, leur servaient à ranger leurs denrées...

... on y distingue également un vieux coffre où ils tiennent leur peu de provision, et une vaste caisse carrée remplie de feuilles de frêne couverte de quelques lambeaux de grosse toile et d’une mauvaise couverture, qui leur sert lieu de grabat sur lequel le vieillard décrépit et l’enfant nouvellement né, la fille et le garçon, le sain et le malade, le mourant et le mort sont souvent confondus....

Leur nourriture ordinaire est de l’orge la plus grossière dont ils font un peu de pain, ou qu’ils font cuire avec le lard le plus rance qui leur sert ensuite de pitance. On les a vus, en dernier lieu, faire le bouillon ou la soupe pour leurs malades avec le vieux oing\footnote{Vieille graisse de porc fondue servant à graisser un mécanisme.}.

Les autres aliments dont ils usent sont les aulx, les oignons et, en hiver, quelques pommes de terre, des raves, les plus mauvais légumes et des fromages encore plus mauvais.

Aussi en est-il mort un très grand nombre.

}

{\fontsize{10.2}{12.0}\selectfont

Dans des temps plus heureux on les a vus se nourrir d’assez bon pain de froment, de bons légumes, d’un peu de viande fraîche, de cochon salé, de harengs et de morue, de châtaignes et d’assez bon fromage et se procurer une quantité suffisante de linge pour se tenir propres, et se coucher sur la paille qu’ils avaient soin de renouveler assez souvent…

Cette maladie, quoique très dangereuse par elle-même n’eut point eu de suites si fâcheuses si, par une fatalité dont on a peu d’exemples, la plupart des malades n’eussent été presque entièrement privés de secours :

Mr Auzilion, chirurgien de Meyrueis, qu’on a eu le malheur de perdre vers le milieu de cette épidémie a été pendant longtemps le seul qui les ait visités et leur ait administré des remèdes. Mais malgré sa bonne volonté et une très grande activité il lui était impossible de les visiter tous et aussi souvent qu’il eut fallu, vu leur nombre, leur éloignement et leur dispersion…

Depuis sa mort jusqu’à mon arrivée, aucun d’eux n’avait été vu que par M. Fignel curé de Lanuéjols.

Aussi en est-il mort un très grand nombre qui n’avaient pris pour tout remède que du suc d’oignon blanc, mêlé avec de la mauvaise huile ou de la suie détrempée dans du vinaigre ; remède réputé souverain parmi eux pour les maladies de la nature de celle-ci, qu’ils qualifiaient de peste blanche…

\medskip
\textit{Pour arrêter les progrès de cette maladie :}

\textbf{1/} On exhortera tous les habitants de la campagne à se tenir aussi proprement qu’il serait possible en leur faisant sentir toute la nécessité d’emporter à des distances assez éloignées le fumier qui entoure leurs chaumières, de nettoyer le dedans de renouveler souvent leurs grabats, de changer de linge et de ne s’exposer au grand air que lorsque les brouillards seront entièrement dissipés.

\textbf{2/} Il sera journellement distribué aux nécessiteux, en danger de tomber malades, à raison de leur détresse, un peu de viande fraîche, du cochon salé, du pain second, quelques légumes et même s’il était possible, un peu de vin et une certaine quantité de linge de rechange.

\textbf{3/} Les médecins feront tous les jours la visite des malades et auront, chemin faisant, la plus grande attention sur ceux qui seraient menacés de la maladie, pour les en garantir, s’il était possible, en leur prescrivant la diète, les boissons les émétique, les purgatifs et les stomachiques convenables.

\textbf{4/} Vu l’état de dépérissement et de faiblesse dans lequel se trouvent la plupart des malades ils soutiendront leurs forces par de bons bouillons, de légers cordiaux et des tisanes légèrement incisives… Ils seront très réservés sur les saignées, les émétique et les purgatifs… En espérant que les approches d’une saison plus tempérée contribuera à terminer bientôt cette fâcheuse maladie.

\vspace{0mm}

\begin{tcolorbox}[
  colback=white,
  colframe=black,
  boxrule=0.7pt,
  arc=0pt,
  left=6pt,right=6pt,top=5pt,bottom=5pt
]
\textbf{On peut penser, sans grand risque d’erreur, que la misère extrême du pays et particulièrement des vallées comme du village de Lanuéjols, autrefois protestant, était une conséquence des persécutions et des dragonnades qui, après le brûlement des Cévennes de 1703, continuèrent tout au long du 18\ieme{} siècle.}
\end{tcolorbox}

}


